\section{Conclusion}
\label{sec:conclusion}

\subsection{Summary}

We have presented a complete algorithmic toolbox for congressional redistricting
that satisfies all eleven legal requirements identified in the post-\callais{}
redistricting landscape.

The toolbox is organised along three orthogonal axes:
Axis A controls what METIS balances (population only, population plus area,
population with VRA alignment, population with partisan vote share);
Axis B controls edge management (geographic boundary lengths, county-sticky
weights, unweighted);
Axis C controls division strategy (balanced bisection, GeoSection, AreaSection,
ApportionRegions, NestSection).

Each legal requirement maps to at most one component on one axis.
The components compose freely except for the Callais mutex: VRASection (Axis~A3)
and ProportionalSection (Axis~A4) are mutually exclusive, enforced at the CLI
level by \texttt{callais\_preflight} and documented in the plan manifest.

The bakeoff across Wisconsin, North Carolina, Georgia, and Alabama demonstrates
five cross-state patterns:
\begin{enumerate}
  \item Geographic voter sorting drives partisan outcomes in all four states;
    the algorithm is not the cause.
  \item VRASection is the only configuration achieving the minority-opportunity
    district counts required by \allenmilligan{} in Alabama and Georgia.
  \item County preservation trades off against compactness at a rate of
    $\sim 4\,\%$ EC cost per $50\,\%$ reduction in county splits.
  \item AreaSection reduces partisan gaps in competitive states by preventing
    urban-peeling geometry, with no systematic partisan direction.
  \item ApportionRegions and NestSection are structurally neutral tools that
    improve auditability and multi-chamber consistency without changing
    partisan outcomes.
\end{enumerate}

\subsection{The Callais Compliance Design}

\callais{}'s disentanglement requirement at p.~36 is the structurally decisive
constraint for algorithmic redistricting.
The toolbox satisfies it by design: each mode uses exactly one signal type,
the two signal-bearing modes are mutex-gated, and every run is documented in
a tamper-detectable plan manifest.

The strong-inference framework of \callais{} holding~2 is directly supported:
a practitioner can run GeoSection (pure political goals) and VRASection (minority
alignment) independently, compare the results, and demonstrate that the political
goals could have been achieved with better minority outcomes.
For Alabama, this comparison confirms the \allenmilligan{} remedy.

\subsection{The Central Empirical Finding}

The central empirical finding of the B-series programme, confirmed across 44
states and three census years, is that
\emph{geographic voter sorting is the dominant driver of partisan redistricting outcomes}.
No redistricting algorithm---however carefully designed---can overcome the
arithmetic of urban-rural political geography.
A state where Democrats are concentrated in three or four major cities will
systematically elect fewer Democratic representatives than a proportional system
would predict, regardless of how the district lines are drawn.

This finding has a constructive implication.
Redistricting litigation based on the claim that ``the algorithm favoured
Republicans'' is not well-founded when the algorithm is the toolbox.
The correct challenge is: ``could the map have achieved the same political goals
with better minority outcomes?'' (the \callais{} strong-inference test).
For the \allenmilligan{} states, the answer is yes---and VRASection provides
the proof.

\subsection{Future Work}

Four extensions are needed to complete the toolbox:

\begin{enumerate}
  \item \textbf{B.11 ApportionRegions}: Implementation and 50-state empirical
    evaluation of the prime-factorisation bisection tree.
  \item \textbf{B.12 ProportionalSection}: Implementation with the
    Callais-compliant mutex gate; empirical comparison against AreaSection
    in state-court proportionality challenge states.
  \item \textbf{B.10 Phase~2--3}: Download and spatial join of TIGER MCD, place,
    and VTD files; parameter sweep at the township and precinct levels for
    MN, WI, and New England states.
  \item \textbf{B.15 Three-Census CSS}: Completion of the 2010 GeoSection sweep
    to compute full three-census CSS for all 50 states.
\end{enumerate}

With these extensions, the toolbox will provide a complete, auditable,
legally-grounded algorithmic framework for the 2031 redistricting cycle.

\paragraph{Data Availability.}
The \texttt{redist} Rust binary implementing GeoSection, VRASection,
AreaSection, NestSection, and StabilitySection is available as open-source
software at \url{https://github.com/apportionment-research/redist}.
All redistricting outputs, adjacency graphs, and analysis scripts are available
at the same repository.
Pre-computed adjacency files for all 50 states are available as GitHub Release assets,
with SHA-256 hashes recorded in the plan manifests.
All parameters are documented in \texttt{docs/REDIST\_CLI.md}.
All results in this paper are reproducible from the stated parameters and the
SHA-256-verified adjacency graphs.
\end{document}
